\bookchapter{Bab 3: Sambungkan Telefon ke “Garis Hayat” — Rangkaian}

\chapterintro{Tanpa rangkaian, telefon pintar kehilangan separuh “kecerdasannya”. Rangkaian adalah garis hayat telefon moden — memahaminya, menggunakannya dengan baik, dan melindunginya adalah mata pelajaran wajib bagi setiap pengguna telefon.}

\booksection{Sifar, Tidak Keliru dengan Rangkaian — Bezakan “Jalan” dan “Kenderaan” Wi-Fi dan Data Mudah Alih}

Bayangkan internet sebagai sebuah bandar.\textbf{“Jalan”}ialah laluan antara telefon anda dan internet,

\begin{tcbinfo}{"Kenderaan"}
ialah data yang bergerak di jalan. Wi-Fi dan data mudah alih adalah dua jalan berbeza, tetapi data (kenderaan) itu sendiri adalah sama.\textbf{1. Perbezaan Asas: Saluran Tempatan vs Lebuhraya Jarak Jauh}Wi-Fi pada asasnya adalah\textbf{teknologi rangkaian kawasan tempatan}, ia seperti jalan tanpa wayar yang dipasang di bawah bangunan anda, bertanggungjawab menghantar data dari penghala ke telefon. Tetapi jika bil lebar jalur lebar belum dibayar atau kabel gentian optik putus, jalan masih ada, tetapi “lebuhraya” di luar telah terputus, semestinya tidak ke mana-mana. Manakala data mudah alih (data mudah alih) dihantar melalui stesen pangkalan 4G/5G pengendali, bersamaan dengan memandu terus ke laluan lebuhraya pengendali, selagi ada isyarat, ia boleh terus menyambung ke internet.

\end{tcbinfo}

\booksubsection{Kos dan Senario: Sewa Bulanan Tetap vs Pengiraan Mengikut Kuantiti}

Wi-Fi biasanya merupakan pelaburan sekali (yuran bulanan jalur lebar), selagi bil elektrik tidak terputus, seluruh keluarga boleh menggunakannya tanpa had, sangat sesuai untuk senario data besar seperti memuat turun filem dan bermain permainan. Manakala data mudah alih dibil mengikut jumlah penggunaan, melebihi pakej mungkin menyebabkan pengurangan kelajuan atau caj tambahan, lebih sesuai untuk tugas ringan dalam senario bergerak seperti di dalam kereta api atau semasa perjalanan kerja.\textbf{Cuba Ini: Semak Pakej Data Anda}Buka [Tetapan] > [Rangkaian Mudah Alih] > [Pengurusan Data] telefon anda, lihat berapa banyak data yang masih tinggal bulan ini, dan App mana yang paling banyak menggunakan data. Tahu keadaan anda, barulah tidak panik apabila keluar.

\booksubsection{Pertimbangan Keselamatan: Wi-Fi Awam adalah “Kolam Pancing” Tersembunyi}

Dari segi keselamatan, data mudah alih biasanya lebih selamat daripada Wi-Fi awam. Data mudah alih disulitkan oleh pengendali, manakala Wi-Fi awam sememangnya mudah diserang oleh penggodam yang membina “titik panas palsu” untuk menangkap data yang tidak disulitkan. Tetapi di sini perlu dibezakan: dengan penggunaan HTTPS secara meluas hari ini, apa yang boleh ditangkap oleh penyerang adalah

\begin{tcbpractice}{laman web mana yang anda layari}
(nama domain), dan

\end{tcbpractice}

\booksubsection{kandungan HTTP yang tidak disulitkan}

; manakala mesej WeChat, kata laluan log masuk App perbankan dan lain-lain yang disulitkan melalui HTTPS, penyerang\textbf{tidak dapat membaca secara langsung}— melainkan melaksanakan serangan penyingkiran SSL yang lebih kompleks dan pengguna mengabaikan amaran sijil. Namun begitu, apabila melakukan operasi pembayaran di luar, menggunakan data mudah alih telefon secara langsung masih merupakan pilihan paling selamat.\textbf{Peringatan Keselamatan}Di tempat awam seperti lapangan terbang, pusat membeli-belah, dan kafe, jangan sambung ke Wi-Fi percuma yang namanya mencurigakan (seperti “Free\_WiFi\_2024”). Walaupun bersambung, jangan sekali-kali log masuk ke App perbankan atau memasukkan kata laluan di atasnya.\textbf{4. Pengesanan di Telefon: Fahami “Kod Rahsia” Bar Status}Bagaimana mengetahui telefon anda kini menggunakan kaedah sambungan apa? Menggunakan telefon Huawei sebagai contoh, apabila

\textbf{atau alat profesional membersihkan sejarah, dan segera tukar kunci yang bocor.}\par

muncul di sebelah isyarat pengendali atau isyarat WLAN dalam bar status, ia bermaksud sedang menggunakan kaedah berkenaan untuk melayari internet; jika kedua-duanya menyala serentak, ini bermakna pecutan rangkaian berkesan, telefon sedang menggunakan kedua-dua saluran secara serentak.

\booksubsection{Pengguna iPhone lebih intuitif: Ikon Wi-Fi menyala bermakna menggunakan Wi-Fi, jika tidak menyala, lihat bar isyarat di penjuru kanan atas bermakna menggunakan data mudah alih. Sebilangan telefon Android juga akan memaparkan perkataan “Hadkan aktiviti latar belakang” atau “4G” dalam bar status, memberitahu anda jenis rangkaian mudah alih semasa secara langsung.}

Satu, Rangkaian Mudah Alih (4G/5G) dan Data Mudah Alih\textbf{Rangkaian mudah alih adalah asas untuk anda menyambung ke internet di mana sahaja. Dari 2G hingga 5G, kelajuan telah meningkat lebih seribu kali ganda, tetapi prinsipnya sebenarnya hampir sama.}Standard Rangkaian:

4G LTE

\booksection{Satu, Rangkaian Mudah Alih (4G/5G) dan Data Mudah Alih}

5G

\booksubsection{: Kelajuan pantas (secara teori tertinggi 10Gbps), tetapi liputan masih dalam penambahbaikan, penggunaan kuasa tinggi. Sesuai untuk memuat turun fail besar dan menonton video definisi tinggi.}

\begin{itemize}
  \item \textbf{5G SA vs NSA}: SA ialah rangkaian mandiri, NSA bergantung pada stesen pangkalan 4G. Pengalaman SA lebih baik, tetapi memerlukan pengendali dan telefon yang menyokong.
  \item \textbf{Hadkan aktiviti latar belakang}Apakah data mudah alih
  \item \textbf{: Ia adalah jumlah data yang dihasilkan semasa telefon melayari internet. Menonton video dan melayari gambar paling banyak menggunakan data, manakala teks dan suara secara relatif menjimatkan.}Tip menjimatkan data
\end{itemize}

\booksubsection{: Aktifkan “Mod Jimat Data” dalam tetapan, kemas kini App atau tonton video definisi tinggi di bawah Wi-Fi.}

\begin{itemize}
  \item \textbf{Bar isyarat $\neq$ kelajuan rangkaian}: Isyarat kuat hanya bermakna sambungan baik, kelajuan rangkaian juga bergantung pada sama ada stesen pangkalan sibuk dan sejauh mana anda dari stesen pangkalan.
  \item \textbf{Amaran Penting: Matikan [Perayauan Data] Sebelum ke Luar Negara!}Menggunakan data mudah alih di luar negara boleh menghasilkan bil yang mahal! Sebelum ke luar negara, pastikan anda mematikan perayauan data dalam tetapan, atau beli kad data tempatan/pakej perayauan antarabangsa terlebih dahulu.
  \item \textbf{Dua, Wi-Fi: Sambungan, Perkongsian dan Keselamatan}Wi-Fi adalah pilihan utama untuk rangkaian di rumah dan pejabat, kelajuan pantas dan tidak menggunakan data mudah alih. Tetapi Wi-Fi juga sentiasa berevolusi, perbezaan antara standard yang berbeza sangat besar.
\end{itemize}

\textbf{Menyambung Wi-Fi:}\par

Pergi ke [Tetapan] > [WLAN] pilih rangkaian dan masukkan kata laluan. Kini banyak telefon menyokong penyambungan melalui imbasan kod, lebih mudah.

\booksection{Dua, Wi-Fi: Sambungan, Perkongsian dan Keselamatan}

Android dan iOS kedua-duanya membolehkan telefon yang telah bersambung berkongsi kata laluan dengan rakan secara pantas. Android menjana kod QR, iOS boleh berkongsi secara automatik apabila berdekatan.

\booksubsection{Evolusi Standard Wi-Fi:}

Standard Wi-Fi ditetapkan oleh IEEE, setiap generasi mempunyai peningkatan yang ketara:

\begin{center}
\begin{tabularx}{\linewidth}{|X|X|X|X|X|}
\hline
\textbf{Standard} & \textbf{Protokol} & \textbf{Jalur Frekuensi} & \textbf{Kadar Maksimum} & \textbf{Ciri} \tabularnewline
\hline
Wi-Fi 4 & 802.11n & 2.4/5GHz & 600Mbps & Standard lama, liputan masih boleh \tabularnewline
Wi-Fi 5 & 802.11ac & 5GHz & 3.5Gbps & Paling meluas digunakan, kelajuan memadai \tabularnewline
Wi-Fi 6 & 802.11ax & 2.4/5GHz & 9.6Gbps & Prestasi berbilang peranti kuat, kependaman rendah \tabularnewline
Wi-Fi 6E & 802.11ax & 2.4/5/6GHz & 9.6Gbps & Menambah jalur 6GHz, tidak sesak \tabularnewline
Wi-Fi 7 & 802.11be & 2.4/5/6GHz & 46Gbps & Pemindahan kelajuan ultra, pengagregatan berbilang pautan \tabularnewline
\hline
\end{tabularx}
\end{center}

\begin{itemize}
  \item \textbf{Teknologi utama Wi-Fi 6}:OFDMA (Akses Berbilang Frekuensi Ortogonal) membolehkan berbilang peranti menghantar data serentak, MU-MIMO (Berbilang Pengguna, Berbilang Input Berbilang Output) meningkatkan kecekapan serentak berbilang peranti, TWT (Masa Bangun Sasaran) menjadikan peranti lebih jimat kuasa.
  \item \textbf{Wi-Fi 6E}:Menambah jalur frekuensi 6GHz berdasarkan Wi-Fi 6, menyediakan lebih banyak saluran, mengurangkan gangguan. Memerlukan telefon dan penghala yang menyokong.
  \item \textbf{Wi-Fi 7}:Memperkenalkan MLO (Multi-Link Operation), boleh guna pelbagai jalur frekuensi serentak untuk pemindahan data, mengurangkan kependaman secara mendadak. Saluran ultra lebar 320MHz membolehkan kelajuan melonjak. Platform cip flagship seperti Snapdragon 8 Elite Gen5, Dimensity 9400 telah menyokong Wi-Fi 7, model seperti Xiaomi 17, Huawei Mate 80, Samsung S26 boleh guna rangkaian Wi-Fi 7 (perlu router Wi-Fi 7).
  \item \textbf{2.4GHz vs 5GHz}:2.4GHz menembusi dinding dengan baik tetapi kelajuannya perlahan dan gangguannya banyak; 5GHz kelajuannya pantas tetapi penembusan dinding lemah. Jalur 6GHz Wi-Fi 6E/7 paling pantas, tetapi penembusan dinding paling lemah, sesuai digunakan dalam bilik yang sama.
\end{itemize}

\booksubsection{Risiko Wi-Fi Awam:}

:Wi-Fi awam mungkin ada penggodam yang menyadap data atau menetapkan hotspot pemancingan. Peraturan:

\begin{itemize}
  \item Elakkan log masuk ke akaun bank, memasukkan kata laluan.
  \item Sebaiknya gunakan penyulitan VPN.
  \item Jangan sambung secara automatik ke Wi-Fi yang tidak dikenali.
  \item Gunakan laman web HTTPS (ikon gembok dalam bar alamat).
\end{itemize}

\begin{itemize}
  \item Tiga, Bluetooth: Pemain Serbaguna Komunikasi Wayarles Jarak Dekat
  \item Bluetooth mungkin merupakan teknologi wayarles yang paling kerap digunakan pada telefon anda — fon telinga wayarles, jam pintar, sistem kereta, papan kekunci dan tetikus, semuanya bergantung padanya.
  \item Prinsip Teknologi:
  \item Bluetooth beroperasi pada jalur ISM 2.4GHz, menggunakan teknologi penyebaran spektrum lompat frekuensi (FHSS) untuk menyebarkan data dan bertukar-tukar penghantaran secara pantas merentasi 79 saluran, mengelakkan gangguan dengan berkesan. Jarak komunikasi Bluetooth dibahagikan kepada tiga kelas: Kelas 1 kira-kira 100 meter, Kelas 2 kira-kira 10 meter (biasa digunakan pada telefon), Kelas 3 kira-kira 1 meter.
\end{itemize}

\booksection{Tiga, Bluetooth: Pemain Serbaguna Komunikasi Wayarles Jarak Dekat}

Bluetooth 4.0

\booksubsection{: Komunikasi antara ahli pasukan semasa mendaki, berkhemah.}

Bluetooth 4.2

\booksubsection{: Meningkatkan kelajuan penghantaran dan perlindungan privasi, menyokong sambungan langsung IPv6.}

\begin{itemize}
  \item \textbf{Bluetooth 5.0}: Kelajuan penghantaran berganda (2Mbps), jarak penghantaran 4 kali ganda (secara teori 300 meter), keupayaan siaran 8 kali ganda. Menyokong audio dwi (menyambung dua fon telinga serentak).
  \item \textbf{Bluetooth 5.1}: Memperkenalkan fungsi pengalamatan arah (AoA/AoD), ketepatan penentuan lokasi boleh mencapai tahap sentimeter.
  \item \textbf{Bluetooth 5.2}: Standard baharu LE Audio, menyokong pengekodan LC3 (kualiti audio lebih baik, penggunaan kuasa lebih rendah), audio berbilang aliran, audio siaran.
  \item \textbf{Bluetooth 5.3/5.4}: Mengoptimumkan lagi penggunaan kuasa dan kestabilan sambungan, mengukuhkan siaran berkala.
  \item \textbf{Penyebaran LE Audio}: Pada 2024-2025, LE Audio telah memasuki fasa penyebaran pantas. Fon telinga benar-benar wayarles baharu jenama utama (seperti Huawei FreeBuds Pro 4, Xiaomi Buds 5 Pro, Samsung Galaxy Buds3 dan lain-lain) telah menyokong LE Audio dan pengekodan LC3, menyediakan kependaman lebih rendah dan kualiti audio lebih baik. Apabila membeli fon telinga baharu, disyorkan untuk mengutamakan produk yang menyokong LE Audio.
  \item \textbf{Senario Penggunaan:}Penghantaran audio
\end{itemize}

\textbf{: Fon telinga wayarles, audio kereta. LE Audio Bluetooth 5.2 menjadikan kualiti audio dan kependaman melonjak secara kualitatif.}: 2025-2026, LE Audio telah memasuki fasa populariti pesat. Fon kepala tanpa wayar baharu jenama utama (seperti Huawei FreeBuds Pro 5, Xiaomi Buds 6 Pro, Samsung Galaxy Buds3 Pro dll) telah menyokong LE Audio dan pengekodan LC3, menyediakan kependaman lebih rendah dan kualiti bunyi lebih baik. Apabila membeli fon kepala baharu, disyorkan utamakan produk yang menyokong LE Audio.

\booksubsection{: Sambungan lebih stabil.}

\begin{itemize}
  \item \textbf{Perpasangan peranti}: Jam pintar, gelang, papan kekunci dan tetikus, pad permainan.
  \item \textbf{Perkhidmatan lokasi}: Bluetooth Beacon digunakan untuk navigasi dalaman, penjejakan barang (seperti Apple AirTag).
  \item \textbf{Rumah pintar}: Rangkaian Bluetooth Mesh mengawal mentol lampu, kunci pintu dan lain-lain.
  \item \textbf{Empat, Kawalan Inframerah: Telefon Berubah Menjadi Kawalan Jauh Serbaguna}Kawalan inframerah adalah salah satu teknologi kawalan wayarles tertua, tetapi masih praktikal pada telefon. Sebahagian telefon Android (seperti beberapa model Xiaomi, Huawei) mempunyai pemancar inframerah terbina dalam, boleh mengawal TV, penyaman udara, kotak atas set dan lain-lain peranti.
  \item \textbf{Rumah pintar}Kawalan inframerah menghantar isyarat denyut cahaya inframerah pada frekuensi tertentu (biasanya 38kHz) melalui diod pemancar cahaya inframerah, dan penerima inframerah peranti menyahkod denyutan ini untuk melaksanakan operasi yang sepadan. Jenama berbeza menggunakan protokol pengekodan berbeza, yang biasa termasuk NEC, RC5, RC6, SONY dan lain-lain. App inframerah dalam telefon mempunyai pustaka pengekodan pelbagai jenama terbina dalam, hanya pilih jenama dan model untuk dipadankan.
\end{itemize}

\booksection{Empat, Kawalan Inframerah: Telefon Berubah Menjadi Kawalan Jauh Serbaguna}

Kearahan

\booksubsection{Menghantar lokasi anda kepada kenalan kecemasan.}

Jarak pendek

\booksubsection{: Jarak berkesan biasanya 5-8 meter.}

\begin{itemize}
  \item \textbf{Apple tidak menyokong}: iPhone tidak mempunyai pemancar inframerah, tidak boleh menggunakan fungsi ini.
  \item \textbf{Lima, NFC: Sihir Sentuhan}NFC (Near Field Communication, Komunikasi Medan Dekat) adalah teknologi komunikasi wayarles frekuensi tinggi pada jarak sangat pendek (dalam 10 sentimeter), beroperasi pada frekuensi 13.56MHz, kadar penghantaran tertinggi 424Kbps.
  \item \textbf{Prinsip dan Mod Komunikasi:}NFC berkembang daripada teknologi RFID, menyokong tiga mod operasi:
\end{itemize}

\booksection{Lima, NFC: Sihir Sentuhan}

: Telefon membaca maklumat dalam tag NFC, seperti membaca pelekat NFC untuk mencetuskan tindakan.

\booksubsection{Mod titik ke titik}

: Dua peranti NFC menghantar data antara satu sama lain, seperti Android Beam (telah dihentikan).

\begin{itemize}
  \item \textbf{Mod simulasi kad}: Telefon menyimulasikan kad pintar, ini adalah mod yang paling kerap digunakan — kad bas, kad akses, kad bank semuanya menggunakan mod ini.
  \item \textbf{Mekanisme Keselamatan:}Jarak ultra pendek
  \item \textbf{: Jarak komunikasi dalam 10 sentimeter itu sendiri adalah jaminan keselamatan fizikal, penyadapan jarak jauh hampir mustahil.}Cip selamat
\end{itemize}

\booksubsection{: Telefon mempunyai cip selamat bebas (SE) terbina dalam atau menggunakan HCE (Simulasi Kad Hos), kunci pembayaran disimpan dalam kawasan selamat, walaupun telefon di-root, ia tidak boleh diekstrak.}

\begin{itemize}
  \item \textbf{Pengesahan cap jari/pengecaman muka}: Pembayaran memerlukan pengesahan biometrik, walaupun telefon hilang, ia tidak boleh digunakan untuk kecurian.
  \item \textbf{Tokenisasi}: Apple Pay, Huawei Pay dan lain-lain tidak menghantar nombor kad sebenar, tetapi menggunakan token sekali guna, walaupun ditangkap, ia tidak boleh digunakan semula.
  \item \textbf{Aplikasi Tipikal:}Kad bas / Kad akses
  \item \textbf{: Menulis maklumat kad fizikal ke dalam telefon, hanya sentuh telefon.}Pembayaran mudah alih
\end{itemize}

\booksubsection{: Apple Pay, Huawei Pay dan lain-lain, mengikat kad bank ke telefon, sentuh untuk membayar.}

\begin{itemize}
  \item \textbf{Perpasangan peranti}: Sentuhan NFC untuk memasangkan fon telinga Bluetooth, pembesar suara dengan pantas.
  \item \textbf{Pengesahan identiti}: Pembacaan kad pengenalan NFC, pengesahan lesen memandu elektronik.
  \item \textbf{Rumah pintar}Apabila anda mempunyai lebihan data, anda boleh mengaktifkan hotspot untuk digunakan oleh tablet dan komputer.
  \item \textbf{Kaedah pengaktifan}: Tetapan > Hotspot Peribadi/Perkongsian Rangkaian, aktifkannya.
\end{itemize}

\booksection{Enam, Hotspot Peribadi: Jadikan Telefon sebagai Penghala Wi-Fi}

: Sangat menggunakan kuasa bateri dan data, ingat untuk menetapkan kata laluan yang kompleks, dan hadkan bilangan peranti yang bersambung.

\begin{itemize}
  \item \textbf{Perkongsian USB}: Gunakan kabel data untuk menyambung ke komputer, lebih stabil dan lebih pantas berbanding hotspot wayarles.
  \item \textbf{Fungsi kunci BL:}Buka [Tetapan], cari "Hotspot Peribadi" atau "Perkongsian Rangkaian", aktifkannya, tetapkan nama dan kata laluan.
  \item \textbf{Gunakan peranti lain (seperti komputer) untuk mengesan dan menyambung ke Wi-Fi ini.}Di pihak telefon, lihat sama ada anda boleh mengurus peranti yang telah bersambung, atau menendang keluar peranti yang tidak dikenali.
\end{itemize}

\begin{tcbpractice}{Bab 4: Pasang dan Urus Aplikasi — "Pakej Kemahiran" Telefon}
\begin{enumerate}
  \item Aplikasi menjadikan telefon mampu melakukan apa sahaja, tetapi perlu diurus dengan baik. Sama seperti anda menyusun tempat duduk dan memperuntukkan kebenaran untuk pekerja baharu, memasang App juga memerlukan satu set proses piawaian.
  \item Satu, Kedai Rasmi: "Supermarket Aplikasi" Paling Selamat
  \item Android
\end{enumerate}

\end{tcbpractice}
